% Resumo — Guia UECE p. 47-49: informativo (finalidade, metodologia,
% resultados, conclusões), voz ativa, 3ª pessoa do singular, 150 a 500
% palavras, parágrafo único. Palavras-chave separadas por ponto e vírgula,
% minúsculas, terminando em ponto.
% Primeira versão (18/11/2026). A skill revisao-resumos-cientificos pede
% resumo "criativo" com menos de 100 palavras quando há texto completo;
% prevalece a norma UECE (150-500 palavras, informativo), na forma do
% resumo convencional da skill. Contagem: ~280 palavras.
% Números: artifacts/reports/results_final_{gpt-oss-120b,qwen3-27b}-tiers.md
% (valores-p com 3 casas, como na Tabela de pares em 05-resultados);
% latência 5,6x / 5,1x como em 05-resultados (sec. de latência).
\setlength{\absparsep}{18pt}
\begin{resumo}
  Agentes baseados em modelos de linguagem geram consultas SQL a partir de
  uma descrição textual do banco de dados, e a correção dessas consultas
  depende do que essa descrição contém. Este trabalho mede o efeito de
  três níveis aditivos de contexto sobre a correção do SQL gerado: a
  definição técnica das tabelas; a mesma definição com descrições de
  catálogo para tabelas e colunas; e o catálogo acrescido de uma camada
  semântica textual, com regras de grão, regras temporais, regras de
  relacionamento, sinônimos e estilo de resposta. O experimento usa uma
  base comercial real descaracterizada, com 11 tabelas em PostgreSQL, 40
  perguntas em português e dois modelos de pesos abertos de famílias
  distintas, \textit{gpt-oss-120b} e \textit{qwen3.8-27b}, acessados por
  uma API gratuita. Cada pergunta é respondida nas três condições, a
  correção é verificada pela igualdade entre o resultado da consulta
  gerada e o de uma consulta de referência, e as condições são comparadas
  pelo teste de McNemar exato com ajuste de Holm. No
  \textit{qwen3.8-27b}, a acurácia passa de 57,5\% para 60,0\% e para
  87,5\%, e a condição semântica supera a de catálogo (\textit{p} ajustado
  de 0,002) e a técnica (\textit{p} ajustado de 0,001). No
  \textit{gpt-oss-120b}, a acurácia passa de 62,5\% para 70,0\% e para
  82,5\%, sem diferença significativa após o ajuste (menor \textit{p}
  ajustado de 0,064). O catálogo não difere da condição técnica em nenhum
  dos modelos. Os erros corrigidos apenas pela camada semântica envolvem o
  tratamento de referências órfãs e a devolução de colunas não pedidas, e
  a latência mediana da condição semântica, medida na camada gratuita do
  provedor, é cerca de 5 vezes a da técnica. Conclui-se que o que aumenta a correção são as regras sobre
  como a base deve ser consultada, e não as descrições de colunas; que um
  arquivo de texto anexado ao \textit{prompt} reproduz o padrão relatado
  pela indústria com plataformas proprietárias, sem custo de licença; e
  que esse contexto atua como instrução, e não como garantia.

  \vspace{\onelineskip}
  \noindent\textbf{Palavras-chave}: text-to-SQL; modelos de linguagem;
  camada semântica; engenharia de \textit{prompt}; catálogo de dados.
\end{resumo}
